\NormalFont\ShortTitle{1 Timóteo}
{\MT Primeira Carta de Paulo a Timóteo

\par }\ChapOne{1}{\PP \VerseOne{1}Paulo, apóstolo de Jesus Cristo segundo o mandado de Deus, nosso Salvador, e do Senhor Jesus Cristo, nossa esperança.
\VS{2}Para Timóteo,
{\ADD{meu}} verdadeiro filho na fé. Graça, misericórdia
{\ADD{e}} paz da parte de Deus, nosso Pai, e de Cristo Jesus, nosso Senhor.
\VS{3}Assim como te roguei quando eu estava indo à Macedônia, fica em Éfeso, para advertires a alguns que não ensinem outra doutrina,
\VS{4}Nem prestem atenção a mitos nem a genealogias intermináveis, que mais produzem discussões que a edificação da parte de Deus na fé.
\VS{5}O fim do mandamento é o amor que vem de um coração puro, e de uma boa consciência, e de uma fé não fingida.
\VS{6}Alguns se desviaram dessas coisas, e se entregaram a discursos inúteis;
\VS{7}querendo ser mestres da Lei, mas não entendendo, nem o que dizem, nem o que afirmam.
\VS{8}Nós, porém, sabemos que a Lei é boa, se alguém usa dela legitimamente;
\VS{9}pois sabemos isto: que a Lei não foi feita para o justo, mas sim, para os injustos e insubordinados, para os ímpios e pecadores, para os irreligiosos e profanos, para os matadores dos pais e das mães, para os homicidas,
\VS{10}para os que cometem pecados sexuais, para os homens que efetuam sexualidade com homens, para os que roubam a liberdade das pessoas,
\FTNT{*}{{\FR 1:10 }
roubam a liberdade das pessoas
Ou: traficantes de escravos
} para os mentirosos, para os que juram falsamente, e qualquer outra coisa contrária à sã doutrina,
\VS{11}conforme o evangelho da glória do Deus bendito,
\FTNT{†}{{\FR 1:11 }
Ou: bem-aventurado
} que foi confiado a mim.
\VS{12}E agradeço àquele que tem me fortalecido, Cristo Jesus, nosso Senhor, porque me considerou fiel, pondo-me no serviço.
\FTNT{‡}{{\FR 1:12 }
Ou: ministério
}
\VS{13}Antes eu era um blasfemo, perseguidor, e opressor; porém recebi misericórdia, pois foi por ignorância que eu agia com incredulidade.
\VS{14}Mas a graça do nosso Senhor Jesus Cristo foi mais abundante, com a fé e amor que há em Cristo Jesus.
\VS{15}Esta é uma palavra fiel e digna de toda aceitação: que Cristo Jesus veio ao mundo para salvar os pecadores, dos quais eu sou o principal.
\VS{16}Mas por isso me foi concedida misericórdia, para que em mim, o principal, Jesus Cristo mostrasse toda a sua paciência como exemplo aos que haviam de crer nele para a vida eterna.
\VS{17}Ao Rei dos tempos, imortal, invisível, o único Deus sábio,{\ADD{seja}} a honra, e a glória, para todo o sempre, amém!
\VS{18}Este mandamento, filho Timóteo, te dou: que, conforme as profecias antes feitas acerca de ti, que batalhes por elas a boa batalha;
\VS{19}mantendo a fé e a boa consciência, que alguns rejeitaram, e naufragaram na fé.
\VS{20}Dentre esses foram Himeneu e Alexandre. Eu os entreguei a Satanás, para aprenderem a não blasfemar.

\par }\Chap{2}{\PP \VerseOne{1}Por isso eu te exorto, antes de tudo, que se façam pedidos, orações, intercessões
{\ADD{e}} atos de gratidão por todas as pessoas;
\VS{2}pelos reis, e por todos os que estão em posição superior, para que tenhamos uma vida quieta e sossegada, em toda devoção divina e honestidade.
\VS{3}Pois isso é bom e agradável diante de Deus nosso Salvador;
\VS{4}que quer que todas as pessoas se salvem e venham ao conhecimento da verdade.
\VS{5}Pois
{\ADD{há}} um só Deus, e um só Mediador entre Deus e os seres humanos: o humano Cristo Jesus.
\VS{6}Ele mesmo se deu como resgate por todos,
{\ADD{como}} testemunho ao seu tempo.
\VS{7}Para isso fui constituído pregador e apóstolo, instrutor dos gentios na fé e na verdade (estou dizendo a verdade em Cristo, não mentindo).
\VS{8}Portanto, quero que os homens orem em todo lugar, levantando mãos santas, sem ira nem briga.
\VS{9}De semelhante maneira, as mulheres se adornem com roupa respeitosa, com pudor e sobriedade, não com tranças, ouro, pérolas, ou roupas caras;
\VS{10}mas sim, com boas obras, como é adequado às mulheres que declaram reverência a Deus.
\VS{11}A mulher aprenda quietamente, com toda submissão.
\VS{12}Porém não permito que a mulher ensine, nem use de autoridade sobre o marido,
\FTNT{*}{{\FR 2:12 }
Ou: homem
} mas sim, que esteja em silêncio.
\VS{13}Pois primeiro foi formado Adão, depois Eva.
\VS{14}E Adão não foi enganado; mas a mulher foi enganada, e caiu em transgressão.
\VS{15}Porém ela se salvará tendo parto,
\FTNT{†}{{\FR 2:15 }
tendo parto
O texto original não é claro se trata do parto de filhos em geral, como em diversas traduções, “dando à luz filhos”, ou se trata apenas do nascimento de Cristo.
} se permanecerem na fé, amor, e santificação com sobriedade.

\par }\Chap{3}{\PP \VerseOne{1}{\ADD{Esta}} palavra
{\ADD{é}} fiel: se alguém deseja ser bispo, deseja uma excelente obra.
\VS{2}É dever, portanto, que o bispo seja irrepreensível, marido de uma mulher, sóbrio, comedido, decente, hospitaleiro, apto para ensinar;
\VS{3}Não dado ao vinho, não violento, não ganancioso, mas sim moderado, não briguento, não cobiçoso por dinheiro;
\VS{4}que governe bem a sua própria casa, e tenha os
{\ADD{seus}} filhos em submissão com toda dignidade.
\VS{5}(Pois, se alguém não sabe governar a sua própria casa, como cuidará da igreja de Deus?)
\VS{6}Não um novo convertido; para que não se torne arrogante, e caia na condenação do diabo.
\FTNT{*}{{\FR 3:6 }
Isto é, da mesma maneira que o diabo foi condenado
}
\VS{7}É dever também que tenha bom testemunho dos que estão de fora, para que não caia em reprovação pública e na cilada do diabo.
\VS{8}De maneira semelhante, os diáconos
\FTNT{†}{{\FR 3:8 }
Ou: servidores
} sejam de boa conduta, não contraditórios, não dados a muito vinho, não cobiçosos de ganho desonesto,
\VS{9}mas mantenham o mistério da fé em consciência pura.
\VS{10}E também esses sejam primeiro testados, depois sirvam, se forem irrepreensíveis.
\VS{11}De semelhante maneira, as mulheres
{\ADD{sejam}} de boa conduta, não maldizentes, sóbrias,
{\ADD{e}} fiéis em todas as coisas.
\VS{12}Os diáconos sejam maridos de uma mulher, e governem bem os filhos e as suas próprias casas.
\VS{13}Pois os que servirem bem obtêm para si uma posição honrosa e muita confiança na fé que há em Cristo Jesus.
\VS{14}Escrevo-te essas coisas esperando ir a ti em breve;
\VS{15}mas, seu eu demorar, para que saibais como se deve comportar
\FTNT{‡}{{\FR 3:15 }
Lit. andar
} na casa de Deus, que é a igreja do Deus vivo, a coluna e o alicerce da verdade.
\VS{16}E, sem dúvida nenhuma, grande é o mistério da devoção divina: Deus foi revelado em carne, justificado em Espírito, visto pelos anjos, pregado aos gentios, crido no mundo, e recebido acima em glória.

\par }\Chap{4}{\PP \VerseOne{1}Mas o Espírito diz expressamente que nos últimos tempos alguns se afastarão da fé, dando atenção a espíritos enganadores, e a doutrinas de demônios,
\VS{2}por meio da hipocrisia de mentirosos, que têm a sua própria consciência cauterizada.
\VS{3}Eles proibirão o casamento, e
{\ADD{mandarão}} se abster dos alimentos que Deus criou para que sejam usados com gratidão pelos fiéis e pelos que conheceram a verdade.
\VS{4}Pois toda criatura de Deus é boa, e não há nada a rejeitar, se for recebida com gratidão.
\VS{5}Porque ela é santificada pela palavra de Deus e
{\ADD{pela}} oração.
\VS{6}Se informares
\FTNT{*}{{\FR 4:6 }
Lit. propuseres
} essas coisas aos irmãos, serás um bom servidor de Jesus Cristo, alimentado com as palavras da fé e de boa doutrina que tens seguido.
\VS{7}Mas rejeita os mitos profanos e de velhas, e exercita-te na devoção divina.
\VS{8}Pois o exercício do corpo tem pouco proveito; mas a devoção divina é proveitosa para tudo, pois tem as promessas da vida atual e da futura.
\VS{9}{\ADD{Esta}} palavra é fiel, e digna de toda aceitação.
\VS{10}Pois para isto também trabalhamos e somos insultados, porque esperamos o Deus vivo, que é o Salvador de todos, especialmente dos crentes.
\VS{11}Manda essas coisas, e ensina.
\VS{12}Ninguém despreze a tua juventude, mas sê exemplo aos crentes, na palavra, no comportamento, no amor, em espírito, na fé, e na pureza.
\VS{13}Persiste na leitura, na exortação e no ensino, até que eu venha.
\VS{14}Não desprezes o dom que está em ti, que te foi dado por profecia, com a imposição das mãos dos presbíteros.
\VS{15}Medita nessas coisas, e nelas te ocupe, para que o teu progresso seja visível por todos.
\VS{16}Tem cuidado de ti mesmo e da doutrina; persevera nessas coisas; porque, se fizeres isso, salvarás tanto a ti mesmo como aos que te ouvem.

\par }\Chap{5}{\PP \VerseOne{1}Não repreendas ao idoso asperamente, mas exorta-o como a um pai; aos jovens, como a irmãos;
\VS{2}às idosas, como a mães; e às moças, como a irmãs, em toda pureza.
\VS{3}Honra as viúvas que são verdadeiramente viúvas.
\VS{4}Mas, se alguma viúva tiver filhos ou netos, aprendam primeiro a exercer piedade com a sua própria família,
\FTNT{*}{{\FR 5:4 }
Lit. casa
} e a recompensar os seus pais; porque isso é bom e agradável diante de Deus.
\VS{5}A que é verdadeiramente viúva e desamparada espera em Deus, e persevera noite e dia em rogos e orações.
\VS{6}Mas a que se entrega aos prazeres, enquanto vive, está morta.
\VS{7}Manda, pois, essas coisas, para que sejam irrepreensíveis.
\VS{8}Porém, se alguém não cuida dos seus, e principalmente dos de sua própria família,
\FTNT{†}{{\FR 5:8 }
Lit. casa
} negou a fé, e é pior que um incrédulo.
\VS{9}A viúva, para ser registrada, não deve ter menos de sessenta anos, e haver sido mulher de um marido.
\VS{10}Ela deve ter testemunho de boas obras: se criou filhos, se foi hospitaleira, se lavou os pés dos santos, se socorreu os aflitos, se seguiu toda boa obra.
\VS{11}Mas não admitas as viúvas jovens; porque, quando têm desejos sensuais contra Cristo, querem se casar;
\VS{12}e têm condenação, por haverem anulado a primeira fé.
\VS{13}E, além disso, também aprendem a andar ociosas de casa em casa; e não somente ociosas, mas também fofoqueiras e curiosas, falando o que não se deve.
\VS{14}Por isso, quero que as mais jovens se casem, gerem filhos, administrem a casa,
{\ADD{e}} não deem nenhuma oportunidade ao adversário de maldizer.
\VS{15}Pois algumas já se desviaram e seguiram Satanás.
\VS{16}Se algum crente ou alguma crente tem viúvas, que as ajude, e não sobrecarregue a igreja, para que ela possa ajudar as que são de verdade viúvas.
\VS{17}Os presbíteros
{\ADD{ou: anciãos}} que lideram bem sejam estimados como dignos de honra dobrada, principalmente os que trabalham na palavra e na doutrina.
\VS{18}Pois a Escritura diz: “Ao boi que debulha não atarás a boca”;
{\RQ{Deuteronômio 25:4
}} e: “o trabalhador é digno do seu salário”.
{\RQ{Lucas 10:7; Levítico 19:13
}}
\VS{19}Não aceites acusação contra um presbítero, a não ser com duas ou três testemunhas.
\VS{20}Aos que pecarem, repreende-os na presença de todos, para que os outros também tenham temor.
\VS{21}Ordeno-te, diante de Deus, do Senhor Jesus Cristo, e dos anjos escolhidos, que guardes essas coisas sem preconceitos, fazendo nada por favoritismo.
\VS{22}A ninguém imponhas as mãos apressadamente, nem participes dos pecados alheios; conserva-te puro.
\VS{23}Não bebas mais
{\ADD{somente}} água, mas usa
{\ADD{também}} de um pouco de vinho, por causa do teu estômago, e das tuas frequentes enfermidades.
\VS{24}Os pecados de algumas pessoas são evidentes antes mesmo do juízo, mas os de algumas são manifestos depois.
\VS{25}Semelhantemente, também, as boas obras são evidentes; e as que não são assim não se podem esconder.

\par }\Chap{6}{\PP \VerseOne{1}Todos os servos que estão sob escravidão
\FTNT{*}{{\FR 6:1 }
Lit. jugo
} estimem aos seus senhores como dignos de toda honra; para que o nome de Deus e a doutrina não sejam blasfemados.
\VS{2}E os que têm senhores crentes não os desprezem por serem irmãos; pelo contrário, que os sirvam melhor, pois os participantes do benefício são crentes e amados. Ensina e exorta estas coisas.
\VS{3}Se alguém ensina doutrina diferente, e não concorda com as sãs palavras do nosso Senhor Jesus Cristo e com a doutrina que é conforme a devoção divina,
\VS{4}é presunçoso, e nada sabe. Em vez disso, se interessa de maneira doentia por questões e disputas de palavras, das quais nascem invejas, brigas, blasfêmias, suspeitas maldosas,
\VS{5}discussões inúteis de pessoas corrompidas de entendimento, e privadas da verdade, que pensam que a devoção divina é um meio de lucrar. Afasta-te dos tais.
\VS{6}Mas grande lucro é a devoção divina acompanhada de contentamento.
\VS{7}Pois nada trouxemos ao mundo, e
{\ADD{é}} evidente que nada podemos levar dele.
\VS{8}Mas se tivermos alimento e algo com que nos cobrirmos, estejamos contentes com isso.
\VS{9}Mas os que querem ser ricos caem em tentação, armadilha, e em muitos desejos insensatos e nocivos, que afundam as pessoas na destruição e perdição.
\VS{10}Pois o amor ao dinheiro é raiz de todos os tipos de males. Alguns o cobiçam, e então se desviaram da fé, e perfuraram a si mesmos com muitas dores.
\VS{11}Porém tu, homem de Deus, foge dessas coisas. Segue a justiça, a devoção divina, a fé, o amor, a paciência, a mansidão.
\VS{12}Que lutes a boa luta da fé; agarra a vida eterna para a qual também foste chamado, e confessaste a boa confissão diante de muitas testemunhas.
\VS{13}Mando-te diante de Deus, que dá vida a todas as coisas, e de Cristo Jesus, que diante de Pôncio Pilatos testemunhou a boa confissão,
\VS{14}que guardes este mandamento sem mancha nem motivo de repreensão, até a aparição do nosso Senhor Jesus Cristo.
\VS{15}Essa nos seus tempos mostrará o bendito
\FTNT{†}{{\FR 6:15 }
Ou: bem-aventurado
} e único Soberano, o Rei dos reis, e Senhor dos senhores;
\VS{16}o único a ter imortalidade, e que habita em luz inacessível; aquele a quem nenhum ser humano viu, nem pode ver; a ele seja honra e poder para sempre! Amém.
\VS{17}Manda aos ricos desta era que não sejam soberbos, nem ponham a esperança na incerteza das riquezas, mas sim no Deus vivo, que nos dá abundantemente todas as coisas
{\ADD{que são}} para alegria;
\VS{18}que façam o bem, sejam ricos em boas obras, dispostos a compartilhar, e generosos.
\VS{19}que acumulem para si mesmos um bom fundamento para o futuro, para que obtenham a vida eterna.
\VS{20}Timóteo, guarda o depósito que
{\ADD{te}} foi confiado, evitando conversas vãs e profanas, e controvérsias do que é falsamente chamado de “conhecimento”;
\VS{21}que alguns declararam seguir, e se desviaram da fé. A graça
{\ADD{seja}} contigo. Amém.
{\ADD{A primeira carta a Timóteo foi escrita de Laodiceia, que é a principal cidade da Frígia Pacaciana}} .\par }